% Generated draft: requires mathematical and source review.
\documentclass[11pt]{article}
\usepackage[margin=25mm]{geometry}
\usepackage{fontspec}
\setmainfont{Latin Modern Roman}
\usepackage{amsmath,amssymb,mathtools,unicode-math}
\setmathfont{Latin Modern Math}
\usepackage{enumitem,needspace,xurl,hyperref}
\setlength{\parindent}{0pt}
\setlength{\parskip}{5pt}
\setlength{\emergencystretch}{3em}
\newcommand{\N}{\mathbb{N}}
\newcommand{\Z}{\mathbb{Z}}
\newcommand{\Q}{\mathbb{Q}}
\newcommand{\R}{\mathbb{R}}
\newcommand{\C}{\mathbb{C}}
\newcommand{\F}{\mathbb{F}}
\newcommand{\A}{\mathbb{A}}
\newcommand{\PP}{\mathbb{P}}
\newcommand{\E}{\mathbb{E}}
\newcommand{\eps}{\varepsilon}
\newcommand{\dd}{\,\mathrm{d}}
\begin{document}
\section*{UnsolvedMath ID: 30000149}
\textbf{Source problem code:} catalogue entry
\subsection{Definitions and mathematical statement}
The question asks for the best attainable accuracy in recovering the edges of an image from noisy projection observations. Let $f$ denote an image, $E(f)$ the edge set to be recovered, and $P_{f,\varepsilon}$ the observation law at noise level $\varepsilon$, with the projection design included in that law. For a specified image class $\mathcal F$ and edge-set loss $\ell$, a standard precise version asks for the minimax risk
\[
 R^*(\varepsilon;\mathcal F,\ell)
 =\inf_{\widehat E}\sup_{f\in\mathcal F}
       \mathbb E_{f,\varepsilon}\bigl[\ell(\widehat E,E(f))\bigr],
\]
where the infimum is over estimators measurable with respect to the noisy projection data. Determining an optimal accuracy rate requires both an attainable upper bound and a matching lower bound in the specified model.

The source does not prescribe image dimension, projection geometry and sampling, noise law, edge regularity or contrast assumptions, or a loss such as boundary-distance or set-measure error. These choices affect the risk and rate. No universal numerical rate, all-angle Radon model or convexity assumption is inserted into the broad question. The displayed minimax notation is the explicit interpretation suggested by the attached assessment, not an assertion that those missing model choices have been supplied.

The literature assessment checked on 21 August 2026 reports that Goldenshluger and Spokoiny obtained a nearly minimax estimator for the standard convex-edge Radon model. It explicitly says that result does not settle the broader extraction. ``Nearly minimax'' is retained without changing it to an exact sharp rate. The partially-solved label is dated triage requiring independent verification.
\subsection{Short English statement}
Determine the optimal statistical accuracy for locating image edges from noisy projections. The supplied convex-edge Radon result is nearly minimax for that model, while the broader question still needs its image class, noise, sampling and loss specified.
\subsection{Sources}
\begin{itemize}
\item \emph{New Inference Concepts for Analysing Complex Data} (2005), Oberwolfach Reports, DOI 10.4171/owr/2004/52, pp. 2795--2838, \url{https://doi.org/10.4171/owr/2004/52}; proposed year supplied as 2004.
\item Literature assessment checked 21 August 2026: original OWR report, \url{https://doi.org/10.4171/owr/2004/52}; Goldenshluger--Spokoiny, \url{https://doi.org/10.1109/TIT.2006.871053}.
\item Local UnsolvedMath registry, permanent ID $30000149$, exact source code \texttt{catalogue entry}, created 25 August 2026 and updated 16 September 2026. Corrected-and-verified status records removal of a redundant problem label.
\end{itemize}
\end{document}
